\chapter{案例库}
\label{ch:cases}

\section{为什么本书需要案例}

如果没有案例库，这个模型就有可能听起来像是放之四海而皆准，
但在操作层面仍然偏薄。案例能够检验同一套语言是否真能解释
不同的行为失败，而不是被牵强地拉伸到失去辨识度。

\section{案例 1：学习启动失败}

\subsection*{情境}
一名学生午饭后打算阅读一个较难的章节，却转而打开短视频信息流，
一刷就是四十分钟。

\subsection*{状态分析}
\begin{itemize}[nosep]
  \item 当前奖励很高，因为新奇感与低努力结合在一起。
  \item 切换成本很高，因为下一段定理证明需要重新装载上下文。
  \item 环境强化很强，因为手机就在手里，书桌也没有提前准备好。
\end{itemize}

\subsection*{解释}
这次失败的主要问题并不是对学习价值缺乏认同，而是一个
障碍几何结构的问题。读者是在没有准备好目标状态的前提下，
进入了午休后的决策窗口。

\subsection*{修复}
在吃饭前先写下下一段证明的具体目标，把相关页面提前翻开，
并把手机物理性地移出工作区。这些步骤会在窗口到来之前压低
$\Switch$ 和 $\Env$。

\section{案例 2：下班后无法开始锻炼}

\subsection*{情境}
一名员工每天早晨都认同规律锻炼的重要性，但在消耗殆尽的工作日之后，
很少真正走进健身房。

\subsection*{状态分析}
\begin{itemize}[nosep]
  \item 到了傍晚，$\Will$ 已经偏低。
  \item 路线不确定、换装摩擦和社交孤立感抬高了 $\DV$。
  \item 一到家，面向休息的环境就会过快地被重新建立。
\end{itemize}

\subsection*{解释}
问题并不是在抽象层面上没有选择健康，而是关键转变恰好发生在
整日循环里最糟糕的节点上。

\subsection*{修复}
把窗口前移：提前装好运动包，选定一套默认训练方案，并让通勤路线
经过健身房。这个干预之所以有效，是因为它把决策从
``我现在要不要临时组织一次锻炼？'' 变成了
``我要不要走进那扇我本来就计划要进的门？''

\section{案例 3：研究工作中的写作回避}

\subsection*{情境}
一名研究生能够读论文、讨论想法，但却反复拖延写作初稿。

\subsection*{状态分析}
\begin{itemize}[nosep]
  \item 奖励梯度偏向研究性浏览，因为这种活动在不暴露模糊表达的情况下
        也会让人感觉自己在产出。
  \item 目标状态带有很大的歧义项：写哪一节，提出什么主张，落成哪一句？
  \item 当前状态保护了自我形象，因为它避免了可见输出。
\end{itemize}

\subsection*{解释}
这是一种伪工作惯性。主体并非没有活动，而是被困在一个邻近状态中；
相较于真正起草，那个状态的短期奖励更安全。

\subsection*{修复}
把下一次写作会话定义为“为第 2 节写出五条粗糙但可见的要点主张”。
这样一来，目标状态就具体到足以进入。质量可以以后再优化；
进入状态这一步不能推迟。

\section{案例 4：睡眠作息崩坏}

\subsection*{情境}
一个人想要更早入睡，却在夜间持续消费高刺激内容，结果睡前时间
不断后移。

\subsection*{状态分析}
\begin{itemize}[nosep]
  \item 夜间内容具有高即时奖励和低阻力。
  \item 目标状态“去睡觉”的行动身份很弱，因为它被当成刺激的缺失，
        而不是一个被设计好的转变。
  \item 环境持续让明亮屏幕和躁动线索保持激活。
\end{itemize}

\subsection*{解释}
主体试图仅凭一条消极规则来打败高奖励状态：``停下来。''
但这条规则并没有给出替代框架。

\subsection*{修复}
设计一个睡前状态，明确地点、物件集合和第一个动作：
调暗灯光、准备纸质书、把充电器放到远离床的位置，并设定固定的关机线索。
这样一来，目标就变成一个积极的状态，而不是空洞的禁止。

\section{跨案例比较}

\begin{center}
\begin{tabularx}{\textwidth}{@{}l l l X@{}}
\toprule
案例 & 主导障碍 & 错过的窗口 & 最佳修复 \\
\midrule
学习启动 & 切换成本 & 休息后的重新进入 & 预装载下一道精确问题 \\
锻炼启动 & 傍晚意志偏低 & 通勤到回家边界 & 调整路线并提前装包 \\
写作回避 & 目标歧义 & 写作开始的第一分钟 & 把目标缩成粗糙但可见的输出 \\
睡眠修复 & 被奖励的默认状态 & 夜间关机转变 & 建立积极的替代仪式 \\
\bottomrule
\end{tabularx}
\end{center}

\section{推荐阅读}

\begin{readingbox}
\textbf{基础}
\begin{itemize}[nosep]
  \item John Holland, \emph{Emergence}，用于理解为什么重复性的局部互动
        会生成稳定的高层模式。
\end{itemize}
\textbf{现代研究}
\begin{itemize}[nosep]
  \item 当这里的案例已经变得具体可感之后，关于 metastability 与
        causal emergence 的实证与综述章节，是最好的延伸阅读。
\end{itemize}
\textbf{前沿}
\begin{itemize}[nosep]
  \item 在更简单的障碍语言尚未成功解释日常案例之前，不要使用前沿的
        consciousness theories 来解释普通案例。
\end{itemize}
\end{readingbox}

\begin{summarybox}
案例库提升了迁移价值。这个模型现在解释的不再只是某一个习惯，
而是数种结构不同的失败。这会让该框架作为教材体系更可信，
也让后续读者获得可复用的模式识别，而不只是孤立轶事。
\end{summarybox}